\chapter{数据说明}

本文实验主要使用 BGL 和 HDFS 两类公开日志数据集。BGL 数据集用于验证系统在大规模服务器日志场景下的窗口级异常检测能力；HDFS 数据集用于辅助比较不同模型在轨迹级日志异常检测任务中的效果。

实验过程中，原始日志首先经过字段抽取、时间排序和文本规范化处理；随后通过动态参数替换生成日志模板，并将模板映射为事件 ID。为了保证训练阶段和测试阶段输入空间一致，检测阶段未见过的模板统一映射为 \texttt{<UNK>} 事件。

数据划分遵循训练集、验证集和测试集相互独立的原则。训练集用于模型参数更新，验证集用于输入序列长度、异常阈值和模型选择，测试集仅用于最终性能评价。本文以 BGL 测试集混淆矩阵作为最终指标计算依据，其中实际正常窗口 \ChFiveData{split.bgl.test.normal} 个，实际异常窗口 \ChFiveData{split.bgl.test.anomaly} 个；模型正确识别正常窗口 \ChFiveData{confusion.bgl.tn} 个，误报 \ChFiveData{confusion.bgl.fp} 个，漏报 \ChFiveData{confusion.bgl.fn} 个，正确识别异常窗口 \ChFiveData{confusion.bgl.tp} 个。由此计算得到 Accuracy 为 \ChFiveData{metrics.bgl.test.accuracy}、Precision 为 \ChFiveData{metrics.bgl.test.precision}、Recall 为 \ChFiveData{metrics.bgl.test.recall}、F1 值为 \ChFiveData{metrics.bgl.test.f1}。
